\subsection{Feedback exposure and concurrent execution}\label{sec:live-study}
The separate forwarding pilot uses identical initial messages, task-bearing tool feedback and randomized session compositions. Six actions require one to three checks, with a controlled subset encountering a recoverable error. \emph{Full} forwards the whole response; \emph{Projected} removes supplementary records while retaining task-required fields. These records deliberately provide a dependency to examine. Projection changes content, encoded length and later execution together.

Two fixed 8B models perform 3,120 local executions over LangGraph, MCP Streamable HTTP and streaming TLS~1.3/HTTP/2. Each model/condition has 156/36/120 train/validation/test tasks in 32/8/24 sessions. Four test sessions are empty. Dedicated and Shared connections permit four workers; Joint adds model-link padding and pacing. Wire checks find overlapping model streams in 107/128 sessions in each shared condition. The local gateway represents a remotely accessed model service on one host. It does not give a network observer access to an in-process model.

\paragraph{First-feedback point differences remain inconclusive}
Full first-feedback LR recognition is 43/120 for Qwen and 40/120 for Llama; Projected gives 23/120 and 22/120. Prefix scores are 25/120 and 19/120. All eight local primary LR/k-FP contrasts exceed .05 after exact Holm correction, including every post-feedback--prefix comparison; the minimum is .0508. Exploratory RF gives 50 versus 30 and 59 versus 19. Complete model-view LR instead changes from 46 to 44 and from 62 to 33; RF changes from 61 to 31 and from 73 to 34. The complete execution retains task fields and model responses, so its exposure is distinct from the first-feedback contrast.

Each local Full/Projected test condition has 518 tool calls, with equal counts in every pair but different call order or actions in 35 Qwen and 26 Llama pairs. The same-count observation excludes fewer calls as the explanation of Llama's decline; it does not isolate feedback content from length and order. The separate hosted-alias retest has 192 attempts and no adjusted primary rejection. It changes network path and backend together. Supplement S1 reports every LR, k-FP, RF and CNN result, all phases, the hosted outcomes and conditional refit intervals.

\paragraph{Connection sharing needs event-level evaluation}
Concurrent supplied-window features pool visible connections during overlapping tasks. Serial-to-concurrent recognition changes therefore combine connection policy and window mixing. Continuous LR/RF detectors instead receive session boundaries and connection identities, without task windows or private stream assignments. Under Shared at IoU .5, instance recall is 35.0--43.3\%, precision 4.3--18.6\%, and false events 8.21--44.79 per test session. At IoU .3, all-active exceeds both learned Shared detectors in F1 for both models; at .5 the ordering reverses. The supplement reports .3/.5/.7 and empty-session controls. These results diagnose the detectors and overlap rule, not general deployment difficulty.

Held-out goal completion is 119/120 for Qwen in all conditions and 117/120 for Llama except Projected at 118/120. Projection costs about .78 of Full's model bytes. Joint shaping costs 19.5 and 22.5 times Shared's model bytes, while model-only LR changes from 25 to 18 and from 19 to 25. These mixed results supply measured design points, with no established end-to-end protection advantage.
